\begin{abstract}
% Role: measured conditional benefit and full-cohort quality limitation.
A GPU panorama-stitching pipeline exposes computational budgets whose cost and
suitability depend on the input pair. This working manuscript studies whether
a scene-aware cross-stage runtime can reduce stitching latency while maintaining
a predefined quality-of-result constraint. The implemented renderer consumes
original RGB images and a global homography, and supports feather blending and
Laplacian multi-band blending with three or five reductions. We define separate
checks for alignment quality, rendering correctness, and composition quality,
together with a same-homography comparison of blending costs. The T4 evaluation
compares a development-selected fixed configuration, a frozen adaptive heuristic,
and an offline per-pair oracle. Of \ChNominal{} nominal test pairs, the fixed
configuration passes relative alignment quality on \ChFixedPass{} and the
heuristic on \ChAdaptivePass{}. The heuristic achieves a \ChSpeedup{}-fold
median paired speedup on the \ChJoint{} common passing pairs. Reduced quality
coverage prevents a full-cohort quality-constrained speedup claim. A separate
known-homography stress test establishes renderer scaling through 4K, while
end-to-end synthetic failures expose registration limits.
\end{abstract}
